\chapter{指揮者のいないマシン}
% full version: chapters/02_glutcomputer.tex

あなたのコンピュータには指揮者がいる。クロック発生器だ。毎秒何十億回も合図を出し、回路の全素子がそろって、足並みをそろえて切り替わる。脳にはそんな指揮者はいない。各ニューロンは自分の都合で働く。信号は好きなときにやって来て、好きなときに出ていく。このようなマシンで、\nt{GLUT}ニューロン、つまり「プラス」だけを使うと何が計算できるかを見てみよう。

\section*{穴のあいたバケツ}

ニューロンは穴のあいたバケツと考えるとわかりやすい。スパイクが1つ届くたびに、ひしゃく1杯の水が入る。水は穴から少しずつ漏れていく。ひしゃくが頻繁に来れば水は縁まで上がり、バケツがひっくり返る。ニューロンがカチッと鳴り、バケツは空になり、また最初からだ。ひしゃくがまれにしか来なければ、水は漏れ出てしまい、何も起こらない。

\pencilfig{2}{\pencilart{2}}{ニューロンは穴のあいたバケツ：入力のしずくが水をため、水は漏れていく。線に届けばカチッ。}

この穴こそが、ニューロンと論理ゲートの最大の違いである。ニューロンは届いたスパイクを永遠には覚えていない。覚えているのは数ミリ秒だけだ。この窓の中に届いたものは足し合わされ、それより前のものは忘れられる。

\section*{「OR」と「AND」}

ひしゃく1杯でバケツがひっくり返るなら、ニューロンはどの入力にも反応する。これが「OR」だ。2杯必要なら、話は面白くなる。指揮者がいない以上、「両方届いたか？」という問いは、\emph{どれだけの時間内に}かを決めない限り意味をなさない。ニューロンが鳴るのは、2つのスパイクがほぼ同時に、最初の1杯が漏れきる前に届いたときだけだ。これが「時間のAND」、つまり同時検出器である。窓は、あるニューロンでは数ミリ秒、聴覚のニューロンでは1ミリ秒未満だ。

\section*{「プラス」だけではできないこと}

「NOT」を作ってみよう。入力が黙っているときに働くニューロンだ。うまくいかない。静寂はバケツに水を注がないからだ。これは1個のニューロンだけでなく、「プラス」だけでできたどんなネットワークにも当てはまる。入力を強めても、そのネットワークのどのニューロンも活動を減らすことはない。ここから最も厄介な性質が出てくる。\emph{活動を活動で消すことはできない}。燃え広がった火を止める方法はただ一つ、薪をくべるのをやめることだけだ。

自然は抜け道を見つけた。多くの感覚器では、同じ刺激を2つのニューロン群が伝える。目では、ある細胞が「明るくなった」と、別の細胞が「暗くなった」と伝える。1本の配線の代わりに「はい」と「いいえ」の配線の組があれば、「NOT」はただで手に入る。配線を入れ替えればいいだけだ。こうした組を使えば、「プラス」だけでどんな論理回路でも組み立てられるし、計算が終わったこともわかる。どの組でもちょうど1本の配線が点灯すれば完了だ。エンジニアはこの方法で、どんな遅延があっても正しく動く非同期回路を作っている。

\section*{遅延から時間を作る}

指揮者がいなければ共通の時計もないが、それでも時間は計算できる。横から来る音は、近い方の耳に遠い方の耳より早く届く。その差は1ミリ秒未満だ。同時検出器の列に沿って、左耳からの配線を左から右へ、右耳からの配線を右から左へ張ってみよう。2つの信号は互いに向かって走り、両方が同時にたどり着く地点で出会う。その場所の検出器が鳴る。時間差が、鳴ったニューロンの\emph{番号}に、つまり音の方向に変わったのだ。精度は10マイクロ秒程度になる。個々のニューロンははるかに不正確なのに、検出器がたくさんあることが救いになる。

\pencilfig{102}{%
% delay line: the left-ear wire runs right along the top, the right-ear wire runs left along the bottom
\draw[pen=pcBlue] (0,0.6) -- (5.6,0.6); \draw[pen=pcBlue] (6.0,-0.6) -- (0.4,-0.6);
\foreach \i in {1,...,7}{
  \pgfmathsetmacro{\x}{0.75*\i}
  \draw[pen=pcBlue] (\x,0.6) -- (\x,0.24); \draw[pen=pcBlue] (\x,-0.6) -- (\x,-0.24);
  \ifnum\i=5 \draw[dotpen=pcAmber, wash=pcAmber, line width=1.1pt] (\x,0) circle (0.2);
  \else \draw[dotpen=pcBlue, wash=pcBlue] (\x,0) circle (0.2); \fi}
\draw[arr=pcBlue] (0.95,0.78) -- (1.75,0.78); \draw[arr=pcBlue] (5.3,-0.78) -- (4.5,-0.78);
\node[lbl, anchor=east] at (-0.05,0.6) {左耳};
\node[lbl, anchor=west] at (6.05,-0.6) {右耳};
\node[lbl, text=pcAmber!70!black, anchor=south] at (3.75,0.95) {ここで出会う};
\draw[arr=pcAmber!70!black] (3.75,0.95) -- (3.75,0.68);
% sound source on the left
\draw[penlite] (-1.25,-0.25) -- (-1.05,-0.25) -- (-0.8,-0.45) -- (-0.8,0.05) -- (-1.05,-0.15) -- (-1.25,-0.15) -- cycle;
\foreach \r in {0.2,0.35}{\draw[penlite] ($(-0.75,-0.2)+(-40:\r)$) arc (-40:40:\r);}
\node[lbl, anchor=north] at (-0.95,-0.5) {左からの音};
}{左からの音は先に左耳に届き、その信号は先へ進む余裕があるので、出会う位置は中央より右になる。鳴ったニューロンの番号が音の方向である。}

\section*{記憶はたき火}

互いに強く興奮させ合うニューロンの群れを考えよう。短い信号で火をつけると、この群れは自分の炎で自分に薪をくべるたき火のように、自らを維持し続ける。信号は終わったのに、群れは燃え続けている。信号があったことを覚えているのだ。これがメモリセルである。信号が短すぎれば、たき火は燃え上がる前に消えてしまう。

だが困ったことがある。この記憶はどうやって消すのか。活動では消せない。ニューロンが疲れて、たき火が燃え尽きるのを待つしかない。

\section*{出来事で起動するタイマー}

群れを鎖のようにつなごう。1つ目が2つ目に火をつけ、2つ目が3つ目に火をつける。1つ目を押すと、ドミノ倒しのように鎖に沿って波が走る。倒れたドミノの番号が、押してからどれだけ時間がたったかを教えてくれる。これはいつも時を刻む時計ではなく、出来事で起動して一度だけ動くタイマーだ。鳴禽が歌っている間、ある脳領域のニューロンは厳密に順番どおりに鳴る。まさにこうした鎖によく似ている。

こうして「プラス」だけでも多くのことができる。足りないのは3つ。多くの中から1つを選ぶこと、命令で記憶を消すこと、そしてネットワークの暴走を防ぐこと。この3つはすべて「マイナス」がもたらす。それが次の章の話である。

\fullversion{2}
